\documentclass[12pt,leqno]{article}

\usepackage{amsmath}
\usepackage{amssymb}
\usepackage{enumitem}
\usepackage{hyperref}
\usepackage{logicproof}
\usepackage{mathtools}

\setlength{\textheight}{9.1in}
\setlength{\topmargin}{-.1in}
\setlength{\headsep}{0in}
\setlength{\oddsidemargin}{-.1in}
\setlength{\textwidth}{6.5in}

\setlength{\parskip}{1em}
\setlength{\parindent}{0pt}

% for logicproof
\setlength{\subproofhorizspace}{.5em}
\setlength{\intersubproofvertspace}{.5em}

% custom symbols
\newcommand{\DefEq}{\stackrel{\text{def}}{=}}
\newcommand{\Mid}{\,\middle\vert\,}

% proof rules
\newcommand{\Intro}[1]{{#1}{\text{I}}}
\newcommand{\IntroA}[1]{{#1}{\text{I}_1}}
\newcommand{\IntroB}[1]{{#1}{\text{I}_2}}
\newcommand{\Elim}[1]{{#1}{\text{E}}}
\newcommand{\ElimA}[1]{{#1}{\text{E}_1}}
\newcommand{\ElimB}[1]{{#1}{\text{E}_2}}

% truth tables
\newcommand{\T}{\mathtt{T}}
\newcommand{\F}{\mathtt{F}}

\title{CS 511 Homework \#4}
\author{Anthony DeRossi}
\date{3 Oct 2024}

\begin{document}
\maketitle

\section*{Exercise 1}

[EML Chapter 1, page 15-16] Exercise 27, part 1

The \textit{$n$-Queens Problem} is the problem of placing $n$ queens on an $n
\times n$ chessboard so that no two queens can attack each other. In this
exercise we specify the requirements of a solution for the $n$-Queens Problem
as a propositional wff $\psi_n$, with one such wff for every $n \ge 4$. For
convenience we use a set $\mathcal{Q}$ of doubly-indexed propositional
variables, instead of $\mathcal{P}$, where the indices range over the positive
integers:

\[
  \mathcal{Q} \DefEq \left\{q_{i, j} \Mid i, j \in
    \left\{1, 2, \ldots\right\}\right\}
\]

The desired wff $\psi_n$ in this exercise is in
$\mathsf{WFF_{PL}}(\mathcal{Q})$. We set the variable $q_{i, j}$ to truth
value $true$ (resp. $false$) if there is (resp. there is not) a queen placed
in position $(i, j)$ of the board, where we take the first index $i$ (resp.
the second index $j$) to range over the vertical axis downward (resp. the
horizontal axis rightward); that is, $i$ is a row number and $j$ is a column
number.

Write the wff $\psi_n$ and justify how it accomplishes its task.

\newpage
\[
  \varphi_n = \varphi_n^{row} \land \varphi_n^{col} \land \varphi_n^{diag1}
    \land \varphi_n^{diag2}
\]

\[
  \varphi_n^{row} = \bigvee \left\{ q_{i, j} \land \bigwedge
    \left\{\lnot q_{i, l} \Mid \begin{array}{c}
      l \ne j \\
      l \in \left\{1, \ldots, n\right\}
    \end{array}\right\} \Mid i, j \in \left\{1, \ldots, n\right\} \right\}
\]

\[
  \varphi_n^{col} = \bigvee \left\{ q_{i, j} \land \bigwedge
    \left\{\lnot q_{k, j} \Mid \begin{array}{c}
      k \ne i \\
      k \in \left\{1, \ldots, n\right\}
    \end{array} \right\} \Mid i, j \in \left\{1, \ldots, n\right\} \right\}
\]

\[
  \varphi_n^{diag1} = \bigvee \left\{ q_{i, j} \land \bigwedge
    \left\{\lnot q_{k, l} \Mid \begin{array}{c}
      i - j = k - l \\
      k \ne i \\
      l \ne j \\
      k, l \in \left\{1, \ldots, n\right\}
    \end{array} \right\} \Mid i, j \in \left\{1, \ldots, n\right\} \right\}
\]

\[
  \varphi_n^{diag2} = \bigvee \left\{ q_{i, j} \land \bigwedge
    \left\{\lnot q_{k, l} \Mid \begin{array}{c}
      i + j = k + l \\
      k \ne i \\
      l \ne j \\
      k, l \in \left\{1, \ldots, n\right\}
    \end{array} \right\} \Mid i, j \in \left\{1, \ldots, n\right\} \right\}
\]

$\varphi_n^{row}$ is satisfied iff there is exactly one queen in each row. If
any propositional variable $q_{i, j}$ is set, then no other variable with the
same $i$ can be set.

$\varphi_n^{col}$ is based on the same principle as $\varphi_n^{row}$, but
with exclusivity on the column $j$ instead of the row.

$\varphi_n^{diag1}$ is satisfied iff there is exactly one queen in each
diagonal. This uses the same exclusivity tactic as above, but with variables
lying along a diagonal instead of a row or column.

$\varphi_n^{diag2}$ is based on the same principle as $\varphi_n^{diag1}$, but
with exclusivity on the antidiagonal.

These four criteria are the requirements for a valid solution. If they are
all satisfied by some assignment of propositional variables, then that
assignment is a solution to the $n$-Queens Problem.

\newpage
\section*{Exercise 2}

[Lecture Slides 13, page 19]

The formal definition of substitution on page 18 can be simplified if every
wff is such that:

\begin{enumerate}
  \item there is at most one binding occurrence for the same variable
  \item a variable cannot have both free and bound occurrences
\end{enumerate}

Formalize this idea.

\medskip
The original wff definition, free variable function, and definition of
substitution are as follows:

\[
  \varphi \Coloneqq \bot \mid \top \mid x \mid (\lnot \varphi) \mid
    (\varphi \land \varphi) \mid (\varphi \lor \varphi) \mid
    (\varphi \to \varphi) \mid (\forall x\ \varphi) \mid (\exists x\ \varphi)
\]

\[
  FV(\varphi) = \left\{
    \begin{array}{ll}
    \varnothing  & \text{if $\varphi = \bot$ or $\top$} \\
    \{x\}        & \text{if $\varphi = x$} \\
    FV(\varphi') & \text{if $\varphi = \lnot\varphi'$} \\
    FV(\varphi_1) \cup FV(\varphi_2) &
      \text{if $\varphi = (\varphi_1 \star \varphi_2)$ and
        $\star \in \{\land, \lor, \to\}$} \\
    FV(\varphi') - \{x\} &
      \text{if $\varphi = (\mathbf{Q}x\ \varphi')$ and
        $\mathbf{Q} \in \{\forall, \exists\}$}
    \end{array}
  \right.
\]

\[
  \varphi[u/x] = \left\{
    \begin{array}{ll}
    \varphi              & \text{if $\varphi = \top$ or $\bot$} \\
    \varphi              & \text{if $\varphi = y$ and $x \ne y$} \\
    u                    & \text{if $\varphi = y$ and $x = y$} \\
    \lnot(\varphi'[u/x]) & \text{if $\varphi = \lnot \varphi'$} \\
    \varphi_1[u/x] \star \varphi_2[u/x] &
      \text{if $\varphi = \varphi_1 \star \varphi_2$ and} \\
      & \text{$\star \in \{\land, \lor, \to\}$} \\
    \mathbf{Q}y\ (\varphi'[u/x])
      & \text{if $\varphi = \mathbf{Q}y\ \varphi'$} \\
      & \text{$\mathbf{Q} \in \{\forall, \exists\}$, $x \ne y$, and} \\
      & \text{$u$ is substitutable for $x$ in $\varphi$} \\
    \varphi & \text{if $\varphi = \mathbf{Q}y\ \varphi'$,} \\
      & \text{$\mathbf{Q} \in \{\forall, \exists\}$, $x = y$}
    \end{array}
  \right.
\]

\newpage
To simplify the substitution definition we can use the following wff
definition:

\[
  \begin{array}{lcll}
  \varphi & \Coloneqq & \bot \mid \top \mid x \mid (\lnot \varphi) & \\
    & \mid & (\varphi_1 \star \varphi_2)
      & \text{if $\star \in \{\land, \lor, \to\}$}, \\
    &&& BV(\varphi_1) \cap BV(\varphi_2) = \varnothing, \\
    &&& BV(\varphi_1) \cap FV(\varphi_2) = \varnothing, \\
    &&& FV(\varphi_1) \cap BV(\varphi_2) = \varnothing \\
    & \mid & (\mathbf{Q}x\ \varphi)
      & \text{if $\mathbf{Q} \in \{\forall, \exists\}$,} \\
    &&& x \notin BV(\varphi)
  \end{array}
\]

Where $FV(\varphi)$ has the same definition as above and the function
$BV(\varphi)$ reflects the bound variables of a wff:

\[
  BV(\varphi) = \left\{
    \begin{array}{ll}
    \varnothing  & \text{if $\varphi = \bot$ or $\top$} \\
    \varnothing  & \text{if $\varphi = x$} \\
    BV(\varphi') & \text{if $\varphi = \lnot\varphi'$} \\
    BV(\varphi_1) \cup BV(\varphi_2) &
      \text{if $\varphi = (\varphi_1 \star \varphi_2)$ and
        $\star \in \{\land, \lor, \to\}$} \\
    BV(\varphi') \cup \{x\} &
      \text{if $\varphi = (\mathbf{Q}x\ \varphi')$ and
        $\mathbf{Q} \in \{\forall, \exists\}$}
    \end{array}
  \right.
\]

This formulation eliminates the ``substitutable'' requirement in the
definition of substitution. Because there is only one binding occurrence for
any given propositional variable, substitution is always valid under a binding
if the substituted variable does not match the bound variable.

\[
  \varphi[u/x] = \left\{
    \begin{array}{ll}
    \varphi              & \text{if $\varphi = \top$ or $\bot$} \\
    \varphi              & \text{if $\varphi = y$ and $x \ne y$} \\
    u                    & \text{if $\varphi = y$ and $x = y$} \\
    \lnot(\varphi'[u/x]) & \text{if $\varphi = \lnot \varphi'$} \\
    \varphi_1[u/x] \star \varphi_2[u/x] &
      \text{if $\varphi = \varphi_1 \star \varphi_2$ and} \\
      & \text{$\star \in \{\land, \lor, \to\}$} \\
    \mathbf{Q}y\ (\varphi'[u/x])
      & \text{if $\varphi = \mathbf{Q}y\ \varphi'$} \\
      & \text{$\mathbf{Q} \in \{\forall, \exists\}$, $x \ne y$} \\
    \varphi & \text{if $\varphi = \mathbf{Q}y\ \varphi'$,} \\
      & \text{$\mathbf{Q} \in \{\forall, \exists\}$, $x = y$}
    \end{array}
  \right.
\]

\newpage
\section*{Problem 1}

[EML Chapter 1, page 15-16] Exercise 27, parts 2, 3, 4

% TODO

\vspace{3\baselineskip}
\hrule

Solutions to the Lean 4 exercises are available on GitHub:

\url{https://github.com/anthonyde/math2001/blob/cs511-fall2024/Math2001/CS511_Homework/hw4.lean}

\section*{Exercise 3}

\section*{Exercise 4}

\section*{Problem 2}

\end{document}
